The proof is in Gauge Theory Gravity. The field is $\hb^{-1}(e_{\tau}) = \gamma_{0}$, $\hb^{-1}(e_{i}) = \sum_{l}E_{li}\gamma_{l}$ with $E = \Delta G\theta_{0}$ (the background labels $\tau, x^{i}$; derived metric $g$ of the statement), so the coframe vectors $f^{a}$ with $\hb^{T}(f^{a}) = \gamma^{a}$ are $f^{0} = e^{\tau}$, $f^{l} = \sum_{i}E_{li}e^{i}$, and $\hb^{T}(e^{i})$ is a combination of the spatial $\gamma^{m}$ only. The rotation-gauge field is the unique solution of the torsion equation read on the constant physical frame, $\hb^{T}(\nabla \wedge f^{a}) = \sum_{b}\gamma^{b} \wedge (\omega(\gamma_{b}) \cdot \gamma^{a})$, $\omega(A) = \Omega(\hb(A))$ (Fact \claimref{synthesis.imported-gtg-dictionary}(ii)); the matrix $\Omega$ of the statement is the coframe's logarithmic rate, a different object. (i) Since $\theta_{0}$ does not depend on $\tau$, $\tau\partial_{\tau}E = (\tau(\partial_{\tau}\Delta)\Delta^{-1} + \Delta\,\Omega\,\Delta^{-1})E = BE$, as $\tau\partial_{\tau}\tau^{p_{l}} = p_{l}\tau^{p_{l}}$; so $\nabla \wedge f^{l} = \tau^{-1}\sum_{m}B_{lm}\,e^{\tau} \wedge f^{m} + (\text{terms } e^{i} \wedge e^{j})$, and
\[
\hb^{T}(\nabla \wedge f^{0}) = 0, \qquad \hb^{T}(\nabla \wedge f^{l}) = \tau^{-1}\sum_{m}B_{lm}\,\gamma^{0} \wedge \gamma^{m} + (\text{spatial bivectors}).
\]
Write $\omega(\gamma_{0}) = \sum_{l}c_{l}\,\gamma_{0} \wedge \gamma_{l} + \rho$ and $\omega(\gamma_{m}) = \sum_{l}k_{lm}\,\gamma_{0} \wedge \gamma_{l} + (\text{spatial})$, $\rho = \sum_{l<m}\rho_{lm}\gamma_{l} \wedge \gamma_{m}$ spatial. From $(\gamma_{0} \wedge \gamma_{n}) \cdot \gamma^{0} = -\gamma_{n}$, $(\gamma_{0} \wedge \gamma_{n}) \cdot \gamma^{l} = \delta^{l}_{n}\gamma_{0}$ and $\rho \cdot \gamma^{l} = \sum_{p}\rho_{pl}\gamma_{p}$: the equation for $a = 0$ reads $-\sum_{n}c_{n}\,\gamma^{0} \wedge \gamma_{n} - \sum_{m,n}k_{nm}\,\gamma^{m} \wedge \gamma_{n} = 0$, so $c = 0$ and $k$ is symmetric; and the $\gamma^{0}$-legs of the equation for $a = l$ (spatial rotation parts of $\omega(\gamma_{m})$ give spatial legs only) read $\tau^{-1}B_{lm} = -k_{lm} + \rho_{lm}$. Hence
\[
\omega(\gamma_{m}) = -\frac{1}{\tau}\sum_{l}\mathrm{sym}(B)_{lm}\,\gamma_{0} \wedge \gamma_{l} + (\text{spatial}), \qquad \omega(\gamma_{0}) = \frac{1}{\tau}\sum_{l<m}\mathrm{antisym}(B)_{lm}\,\gamma_{l} \wedge \gamma_{m},
\]
the boost block and the $e^{0}$-leg of the rotation block, exactly and also for $x$-dependent $\theta_{0}$; as Cartan one-forms (the right-action values, the negative of $\Omega$ on vectors) the boost block is $\omega^{l0} = \tau^{-1}\mathrm{sym}(B)_{lm}e^{m}$. For $G = 1$, $B = P$ and $\omega(\gamma_{m}) = -(p_{m}/\tau)\gamma_{0} \wedge \gamma_{m}$, the boost block of Lemma \claimref{dark-energy.variable-exponents-residue-pointwise-corrections-are}. The second fundamental form of the slices is the map $A \mapsto (A \cdot D)\gamma_{0}$ on physical slice vectors ($\gamma_{0}$ is the unit normal, constant in the physical frame), and $(\gamma_{m} \cdot D)\gamma_{0} = \omega(\gamma_{m}) \cdot \gamma_{0} = \tau^{-1}\sum_{l}\mathrm{sym}(B)_{lm}\gamma_{l}$: in the orthonormal physical frame $\tau K$ is $\mathrm{sym}(B)$, and the mixed coordinate form $K^{i}{}_{j}$ is its conjugate by $E$. The check solves the torsion condition exactly with the engine at $\tau_{0} = 1$ and $\tau_{0} = 2^{7}$ (where $\tau^{p_{l}}$ is rational) for two shears $\exp(c\log\tau\,N)$, one in each order, and the shear $1 + \tau^{3/7}E_{12}$, with an $x$-dependent $\theta_{0}$, and confirms the two displayed blocks and $(\gamma_{m} \cdot D)\gamma_{0}$ entry by entry. \emph{In components} (a reading): $\partial_{\tau}e = \tau^{-1}Be$ gives $de^{l} = \tau^{-1}B_{lm}\,e^{0} \wedge e^{m} + d_{x}e^{l}$, matched in Cartan's first structure equation by the symmetric boost block and the antisymmetric $e^{0}$-leg; and $K = \tfrac12 h^{-1}\partial_{\tau}h$ for $h = e^{\mathsf T}e$ is conjugate to $\tau^{-1}\mathrm{sym}(B)$, which the earlier check computes from the coordinate metric at the Kasner point $p = (-\tfrac27, \tfrac37, \tfrac67)$ (twelve cases; for $G = 1$ the characteristic polynomial of $\tau K$ is that of $P$, $(1, 0, -\tfrac{36}{343})$).

(ii) $P$ is constant, so $\tau^{-1}\mathrm{sym}(B)$ and $\tau^{-1}\mathrm{antisym}(B)$ are $O(\tau^{-1})$ if and only if $B$, equivalently $\Delta\Omega\Delta^{-1}$, is bounded, and their residues are $P$ and $0$ if and only if $\Delta\Omega\Delta^{-1} \to 0$, whose entries are $\Omega_{lm}\tau^{p_{l} - p_{m}}$. If $\Omega$ is bounded and $\Omega_{lm} \neq 0$ only where $p_{l} > p_{m}$, every nonzero off-diagonal entry carries a positive power of $\tau$ and tends to zero, and the diagonal of a strictly triangular $\Omega$ vanishes, so the residue is unchanged. For $G = \exp(c\log\tau\,N)$, $N$ commutes with $G$, so $\Omega = \tau\cdot c\tau^{-1}N = cN$; for $N$ strictly triangular in some order each unordered pair $\{l, m\}$ carries at most one of $N_{lm}$, $N_{ml}$, the diagonal of $\mathrm{sym}(B)$ is $P$, its $(l,m)$ and $(m,l)$ entries are $\tfrac{c}{2}N_{lm}\tau^{p_{l} - p_{m}}$, and $\mathrm{tr}((\tau K)^{2}) = \mathrm{tr}(\mathrm{sym}(B)^{2})$ is the displayed sum. A term with $p_{l} < p_{m}$ grows without bound, and $\mathrm{tr}((\tau K)^{2}) = \tau^{2}K^{i}{}_{j}K^{j}{}_{i}$ is a scalar, so no change of frame removes the growth. The script pins the sums at $t = \tfrac12, \tfrac14, \tfrac18, \tfrac1{16}$ for one $N$ of each kind ($\tfrac{821657}{819200}$ decreasing to $1$, and $\tfrac{1184569}{25}$ increasing).

(iii) For $G = 1 + fE_{12}$ with $f = \tau^{3/7}$, $G^{-1} = 1 - fE_{12}$ and $\Omega = \tau f'E_{12} = \tfrac37\tau^{3/7}E_{12}$, which tends to zero (the frame converges to the identity at rate $f' = O(\tau^{-4/7}) = o(1/\tau)$), while $(\Delta\Omega\Delta^{-1})_{12} = \tfrac37\tau^{3/7 + p_{1} - p_{2}} = \tfrac37\tau^{-2/7}$ is unbounded, $p_{1} < p_{2}$. By (ii), $\mathrm{tr}((\tau K)^{2}) = 1 + 2\bigl(\tfrac{3}{14}\tau^{-2/7}\bigr)^{2} = 1 + \tfrac{9}{98}\tau^{-4/7}$; the script computes it from the coordinate metric at four values of $t$ ($\tfrac{121}{49}, \tfrac{1201}{49}, \tfrac{18481}{49}, \tfrac{294961}{49}$).

(iv) In GTG, with the unit normal $\gamma_{0}$ and $K = \tau^{-1}\mathrm{sym}(B)$ the second fundamental form of (i), the store-sign Ricci map and Einstein map along $\gamma_{0}$ are
\[
\gamma_{0} \cdot \mathcal{R}(\gamma_{0}) = \partial_{\tau}\mathrm{tr}K + \mathrm{tr}(K^{2}), \qquad -2\,\gamma_{0} \cdot \mathcal{G}(\gamma_{0}) = (\mathrm{tr}K)^{2} - \mathrm{tr}(K^{2}) + s\,R^{(3)},
\]
with $R^{(3)}$ the scalar curvature of the slices and $s = \pm 1$ its sign convention (the standard Raychaudhuri and Gauss identities for a synchronous field, cited from memory; the check confirms both displays exactly, on flat slices, for the non-diagonal fields of (i) with constant $\theta_{0}$, by computing the field strength of the engine's rotation-gauge field). The Einstein condition $\mathcal{R}(a) = \Lambda a$ gives $\gamma_{0} \cdot \mathcal{R}(\gamma_{0}) = \Lambda$ and $-2\gamma_{0} \cdot \mathcal{G}(\gamma_{0}) = 2\Lambda$. With $\tau K = R + o(1)$ and $\tau\partial_{\tau}(\tau K) \to 0$, $\partial_{\tau}\mathrm{tr}K = -\tau^{-2}\mathrm{tr}R + o(\tau^{-2})$, so the first gives $\mathrm{tr}R^{2} - \mathrm{tr}R = O(\tau^{2})$, hence $\mathrm{tr}R = \mathrm{tr}R^{2}$; with $\tau^{2}R^{(3)} \to 0$ the second gives $(\mathrm{tr}R)^{2} = \mathrm{tr}R^{2}$. (In components the first is $R_{00} = -\partial_{\tau}\mathrm{tr}K - \mathrm{tr}(K^{2})$ in the round-sphere-positive sign, $\mathrm{Ric}^{\circ} = -\mathrm{Ric}$.) So $\mathrm{tr}R \in \{0, 1\}$, and for $\mathrm{tr}R > 0$ the real eigenvalues of the symmetric $R$ satisfy $\sum = \sum(\cdot)^{2} = 1$, the Kasner circle; $\Lambda$ enters only at $O(\tau^{2})$, as in Proposition \claimref{dark-energy.cosmological-constant-does-not-change-residue}. The relations are the leading field equations of Proposition \claimref{dark-energy.frozen-point-datum-kasner-wall}(ii) read for a non-diagonal residue, and under the criterion of (ii) $R = P$. The script checks the residue algebra at the chosen point and a control off the circle. The criterion of (ii) is the symmetry-wall structure of the Iwasawa-frame analysis of cosmological billiards (Damour, Henneaux and Nicolai 2003; Heinzle, Uggla and Rohr 2009; cited from memory); the computation here is self-contained. The hypothesis $\tau^{2}R^{(3)} \to 0$ is not derived here, and it is not generic: on the Kasner circle some exponent is $\leq 0$, and for generic vacuum data with a negative exponent the slices' curvature is expected to dominate and produce a bounce (argued from the same literature; filed alongside, synthesis.unit-boost-residue-needs-nonpositive-exponent); in the diagonal class the momentum constraints of \claimref{synthesis.diagonal-kasner-momentum-constraint} are the next order, and their non-diagonal analogue is not treated.
